\documentclass[12pt,a4paper]{ctexart}

% ---------- 基础包 ----------
\usepackage[a4paper,margin=2.4cm]{geometry}
\usepackage{amsmath,amssymb,amsthm,mathtools,bm}
\usepackage{mathrsfs}
\usepackage{enumitem}
\usepackage{booktabs}
\usepackage{hyperref}
\usepackage[nameinlink,capitalise,noabbrev]{cleveref}
\usepackage{tikz-cd}
\usepackage{xcolor}

% ---------- 链接样式 ----------
\hypersetup{
  colorlinks=true,
  linkcolor=blue!60!black,
  citecolor=teal!50!black,
  urlcolor=purple!60!black,
  pdftitle={如何重新发明Beout定理？},
  pdfauthor={Libin Wang},
}

% ---------- 段落与列表 ----------
\setlength{\parindent}{2em}
\setlength{\parskip}{0.3em}
\setlist[itemize]{itemsep=0.25em,topsep=0.2em}
\setlist[enumerate]{itemsep=0.25em,topsep=0.2em}

% ---------- 定理环境 ----------
\theoremstyle{plain}
\newtheorem{theorem}{定理}[section]
\newtheorem{proposition}[theorem]{命题}
\newtheorem{lemma}[theorem]{引理}
\newtheorem{corollary}[theorem]{推论}
\newtheorem{conjecture}[theorem]{猜想}

\theoremstyle{definition}
\newtheorem{definition}[theorem]{定义}
\newtheorem{example}[theorem]{例}
\newtheorem{exercise}[theorem]{习题}
\newtheorem{problem}[theorem]{问题}

\theoremstyle{remark}
\newtheorem*{remark}{注}
\newtheorem*{idea}{思路}

% ---------- 常用数学命令 ----------
\newcommand{\N}{\mathbb{N}}
\newcommand{\Z}{\mathbb{Z}}
\newcommand{\Q}{\mathbb{Q}}
\newcommand{\R}{\mathbb{R}}
\newcommand{\C}{\mathbb{C}}
\newcommand{\F}{\mathbb{F}}
\newcommand{\G}{\mathbb{G}}
\newcommand{\fp}{\mathbb{F}_p}

\newcommand{\abs}[1]{\left\lvert #1 \right\rvert}
\newcommand{\norm}[1]{\left\lVert #1 \right\rVert}
\newcommand{\set}[1]{\left\{ #1 \right\}}
\newcommand{\ang}[1]{\left\langle #1 \right\rangle}
\newcommand{\ideal}[1]{\left( #1 \right)}

\DeclareMathOperator{\ord}{ord}
\DeclareMathOperator{\lcm}{lcm}
\DeclareMathOperator{\Hom}{Hom}
\DeclareMathOperator{\End}{End}
\DeclareMathOperator{\Aut}{Aut}
\DeclareMathOperator{\Gal}{Gal}
\DeclareMathOperator{\im}{im}
\DeclareMathOperator{\Ker}{Ker}
\DeclareMathOperator{\charac}{char}

% ---------- 学习笔记信息 ----------
\newcommand{\CourseTitle}{为什么乘法逆元唯一？}
\newcommand{\AuthorName}{王立斌}
\newcommand{\Term}{2026 春季}

\title{\vspace{-1.2cm}\textbf{\CourseTitle}}
\author{\AuthorName}
\date{\Term}

\begin{document}
\maketitle
\tableofcontents
\vspace{0.8em}
\hrule
\vspace{1em}

\section{动机}
欧几里得算法引出了B\'zout定理，B\'zout定理之后，我们学习了扩展欧几里得算法。并且进一步引入了我们本课程的黄金定律：\emph{消去律}。即：
\begin{proposition}[消去律]
\label{prop:cancel}
如果$\gcd(c, m) = 1$，则对于任意整数$a$和$b$，如果$ac \equiv bc \pmod{m}$，则有$a \equiv b \pmod{m}$。
\end{proposition}

消去律从表面上看是使用除法，把等式两边共有的因子$c$除掉。然而，我们引入\emph{乘法逆元}的概念，即：
\begin{definition}
\label{def:inverse}
如果存在$c^{-1}\in \mathbb{Z}$，满足$c c^{-1} \equiv 1 \pmod{m}$。称$c^{-1}$为$c$模$m$的乘法逆元。
\end{definition}
那么，实际上消去律做的不是除法，而是乘法，即在等式$ac \equiv bc \pmod{m}$两边乘以$c^{-1}$，得到$a \equiv b \pmod{m}$。

乘法逆元源自于B\'ezout定理,求乘法逆元的算法本质上是扩展欧几里得算法。我们还声称，如果$\gcd(c, m) = 1$，则$c$模$m$的乘法逆元存在且唯一。现在的问题是：B\'ezout系数是否唯一？如果不唯一，为什么乘法逆元会唯一？为此，我们需要重新审视B\'ezout系数的形式。

\section{B\'ezout系数不唯一}
回顾B\'ezout定理，并考察B\'ezout系数的形式。
\begin{theorem}[B\'ezout定理]
\label{thm:bezout}
对整数 $a,b$，存在整数 $r,s$ 使得
\[ar+bs=\gcd(a,b).
\]\end{theorem}


考虑互素的两个整数$a$和$b$，根据B\'ezout 定理，则存在$r, s \in \mathbb{Z}$使得：
	\[ar + b s = 1\]
从以上等式出发，现在想构造出新的B\'ezout系数应该怎么做呢？

一种初始的想法：找一个未知量$e$
		\[ a r + bs + e - e = 1\]
然后分别把$e$与$ar$和$bs$“结合”起来：
\begin{equation}
\label{eq:init}
  ar + e + bs -e = 1
\end{equation}


此时，我们希望$e$与$ar$(或者$bs$)有共同的因子使得：
		\[ a(r + e') + b(s - e') = 1\]
请问，此时，我们判断$e$应该具有什么形式？

要使得$e$中能分别抽取出$a$和$b$，难道不是$e$必然同时具有$a$和$b$这两个分量吗？即$e = kab$，$k$是任意整数。把$e$带入方程\ref{eq:init}，通过简单计算得：
		\[ a r + kab + bs - kab = a(r +kb) + b(s - ka) = 1\]
也就是说，我们构造出了一系列的B\'ezout系数$r + kb$和$s -k a$。B\'ezout系数不唯一。

\section{乘法逆元唯一性的分析}
给定$a$和$b$，$r + kb$和$s -k a$确实是B\'ezout系数。
但是，所有的B\'ezout系数是否都形如$r + kb$和$s -k a$呢？为此，我们需要更多的分析。

假设同时存在两组B\'ezout系数满足：
\begin{eqnarray*}
  a r_0 + b s_0 = 1\\
  a r_1 + b s_1 = 1
\end{eqnarray*}

请问，此时，我们如何把$r_0$(或者$s_0$)表达成$r_1$与$b$（或者$s_1$与$a$）组成的形式？请回忆中学的“消元法”！


这种中学生应该熟悉的思路是：消去两个等式中含有$b$(或者$a$)的部分。即式子一左右两边同乘$s_1$，式子二左右同乘$s_0$，所得相减：
		\[ a r_0 s_1 - a r_1 s0 = a(r_0 s_1 - r_1 s_0) = s_1 - s_0\]
同样的操作：
		\[b s_0 r_1 - b s_1 r_0 = b(s_0 r_1 - s_1 r_0) = r_1 - r_0\]

		令$k = r_0 s_1 - r_1 s_0$，可得：
		\[r_0 = r_1 + kb\]
		\[s_0 = s_1 - ka \]
  也就是说，任意两组B\'ezout系数$r_0, s_0$和$r_1, s_1$之间的关系是$r_0 = r_1 + kb$和$s_0 = s_1 - ka$。这也说明，所有的B\'ezout系数都形如$r + kb$和$s -k a$。

  回到乘法逆元的唯一性问题，$c$模$m$的乘法逆元是B\'ezout系数中的$r \bmod m$，而且$r$的形式是$r + km$。虽然$r$不唯一，但是$r$在模$m$意义下的值都相同。因此$c$模$m$的乘法逆元唯一。

\section{小结}
请注意，以上分析不是仅仅为了证明乘法逆元的唯一性。如果仅仅为了证明乘法逆元的唯一性，完全可以直接构造以下证明。
\begin{proof}[$c$模$m$的乘法逆元唯一.]
如果存在两个不同的整数$c_0$和$c_1$是$c$模$m$的乘法逆元，则有$cc_0 \equiv 1 \pmod{m}$和$cc_1 \equiv 1 \pmod{m}$，两式相减得到$cc_0 - cc_1 \equiv 0 \pmod{m}$，即$c(c_0 - c_1) \equiv 0 \pmod{m}$。由于$\gcd(c, m) = 1$，$m$不能整除$c$，所以$m$整除$(c_0 - c_1)$，即$c_0 \equiv c_1 \pmod{m}$，即乘法逆元在模$m$意义下唯一。
\end{proof}

因此，本文的分析重点在回归到B\'ezout定理，分析B\'ezout系数的形式，进而分析乘法逆元的唯一性。从而得到更清晰的知识概貌。

以上分析受到了FINT\cite{fint}的启发。记录下来的意图并非为了展示新的知识点，而为了强调：回归到中学知识，使用最基本的方法，无需任何高级的技巧，就可以分析理解新的知识点。从目前大家的学习状况看，缺乏阅读，缺乏思考是最大的问题。

希望这篇小作文对大家有所启发(20260321)。

\begin{thebibliography}{9}

\bibitem{fint}
Joseph H. Silverman.
\newblock \emph{A Friendly Introduction to Number Theory}.
\newblock 4th ed.
\newblock Pearson, 2012.

\end{thebibliography}
\end{document}
